The study will assemble a stratified corpus of Mandarin-language popular recordings released between 2000 and 2020 with approximately even sampling by year and documented inclusion/exclusion rules; each track will be assigned a persistent identifier and linked metadata (release year, circulation channel/provenance where recoverable, artist identity, credited roles such as vocalist/composer/producer where available, and artist gender coding with documented uncertainty). Audio will be standardized prior to analysis, and mid-level musical representations will be derived via candidate MIR pipelines: predominant melody could be estimated using contour-based approaches \citep{salamon2012melody} to compute pitch range, interval distributions, contour archetypes, and repetition statistics; beat and downbeat grids could be estimated using widely adopted tracking pipelines (e.g., implemented in \textit{madmom}) \citep{bock2016madmom}; metric interpretations could be encoded probabilistically where analytic ambiguity is plausible \citep{declercq2016measuring}. Formal organization could be modeled through structural segmentation based on self-similarity and repetition, using spectral clustering approaches \citep{mcfee2014structure,nieto2020structure}; segment boundaries and recurrence patterns could be translated into theory-facing representations (e.g., verse/chorus/bridge) using repetition strength, timbral stability, and (where feasible) lyric alignment, with limited expert correction on a validation subset. Harmonic organization could be represented using chroma-based descriptors and coarse loop typologies, supplemented where reliable by automated chord estimates benchmarked against expert-annotated resources \citep{burgoyne2011chord}. Production and timbral characteristics could be modeled using spectral features, loudness proxies, and embedding-space representations (e.g., computed with \textit{librosa}) \citep{mcfee2015librosa}, informed by perceptual approaches to timbre dimensions \citep{McAdams1995Perceptual}. Formal interpretation will draw selectively on popular-music form scholarship \citep{covach2005form,nobile2020} and on functional concepts adapted from the \textit{Formenlehre} tradition \citep{caplin1998,Schmalfeldt2009Process,VandeMoortele2015Formal}. Homogeneity will be operationalized using within-year dispersion measures and between-year distributional distances across formal, melodic, harmonic, and timbral features, including prototype concentration/entropy for structural patterns; change-point analyses will compare hypothesized break years (2009, 2012, 2018) with data-driven alternatives, with uncertainty assessed via resampling; gender- and role-stratified analyses will evaluate whether trajectories differ by artist gender and credited roles where available. An LLM-based agent workflow will generate and organize candidate semantic relationships linking quantitative outputs (feature trends, prototype shifts, change points, robustness results) to interpretive constructs related to platformization and gender politics, with candidate relationships recorded as a claim ledger (inputs, prompts, outputs, and expert decisions) for traceability. Approximately 10\% of the corpus will be manually audited to assess extraction and segmentation reliability, and longitudinal claims will be accompanied by robustness checks under alternative feature sets, parameterizations, and corpus constructions.